% defer/rcuAPI.tex
% mainfile: ../perfbook.tex
% SPDX-License-Identifier: CC-BY-SA-3.0

\subsection{RCU Linux内核API}
\label{sec:defer:RCU Linux-Kernel API}
\OriginallyPublished{sec:defer:RCU Linux-Kernel API}{RCU Linux-Kernel API}{Linux Weekly News}{PaulEMcKenney2008WhatIsRCUAPI}

本节从Linux内核API的角度来审视RCU。\footnote{
	用户空间RCU的API在其他文献中有文档说明~\cite{PaulMcKenney2013LWNURCU}。}
\Cref{sec:defer:RCU API and Software Engineering}
探讨了RCU的API与软件工程之间的关系，
\cref{sec:defer:RCU has a Family of Wait-to-Finish APIs}
介绍了RCU的等待完成API，
\cref{sec:defer:RCU has Publish-Subscribe and Version-Maintenance APIs}
介绍了RCU的发布-订阅和版本维护API，
\cref{sec:defer:RCU has List-Processing APIs}
介绍了RCU的链表处理API，
\cref{sec:defer:RCU Has Diagnostic APIs}
介绍了RCU的诊断API，
\cref{sec:defer:Where Can RCU's APIs Be Used?}
描述了RCU的各种API可以在哪些上下文中使用。
最后，
\cref{sec:defer:So; What is RCU Really?}
给出了结论性评述。

对内核内部不太感兴趣的读者可以跳到
\cpageref{sec:defer:RCU Usage}上的
\cref{sec:defer:RCU Usage}，
但最好在阅读完下一节关于软件工程考量的内容之后再跳。

\subsubsection{RCU API与软件工程}
\label{sec:defer:RCU API and Software Engineering}

已经提前查看过
\cref{tab:defer:RCU Wait-to-Finish APIs,%
tab:defer:RCU Markers for Read-Side Critical Sections,%
tab:defer:SRCU Markers for Read-Side Critical Sections,%
tab:defer:RCU Polled Update-Side Primitives,%
tab:defer:RCU Publish-Subscribe and Version Maintenance APIs,%
tab:defer:RCU-Protected List APIs,%
tab:defer:RCU Diagnostic APIs}
的读者可能已经注意到，Linux内核API的完整列表
拥有100多个成员。
这与\cref{tab:defer:Core RCU API}中仅显示的六个
API成员形成了鲜明的（也许令人沮丧的）对比。
这种情况显然引出了``为什么这么多？？？''这个问题。

以下各节将更彻底地回答这个问题，
但在此期间本节其余部分概述了这些动机。

有一句古老的智慧："人非圣贤，孰能无过。"
这意味着RCU API中相当大一部分的目的是
提供诊断功能，最显著的是\cref{tab:defer:RCU Diagnostic APIs}中的内容，
但其他地方也有。

人为错误的重要原因之一是人脑的局限性，
例如短期记忆的有限容量。
本书中展示的玩具示例不会触及这些极限。
这是出于必要：
许多读者在学习新材料时会达到认知极限，
因此示例需要保持简单。

因此，这些示例将\co{rcu_dereference()}调用
保持在与外围\co{rcu_read_lock()}和
\co{rcu_read_unlock()}调用相同的函数中。
相比之下，现实世界的软件经常需要从
不同的函数甚至不同的翻译单元中调用这些API成员。
Linux内核RCU API因此已扩展以适应lockdep，
lockdep允许\co{rcu_dereference()}及其相关函数在未受
\co{rcu_read_lock()}保护时发出警告。
Linux内核RCU还检查一些双重释放错误、
RCU read-side critical section中的无限循环，
以及在RCU read-side critical section中尝试调用
quiescent state的情况。

现实世界软件适应人类认知局限性的另一种方式是
通过抽象。
因此，Linux内核API除了
\cref{tab:defer:Core RCU API}中面向指针的核心API之外，
还包括操作链表的成员。
Linux内核本身也提供了受RCU保护的哈希表和搜索树。

像Linux这样的操作系统内核运行在
\cref{fig:intro:Software Layers and Performance; Productivity; and Generality}
所示的软件栈"铁三角"的底部附近，
性能至关重要。
因此，许多RCU API都有专用于快速路径的特化变体，
例如，如
\cref{sec:defer:RCU has Publish-Subscribe and Version-Maintenance APIs}所讨论的，
\co{RCU_INIT_POINTER()}可以在受RCU保护的指针
被赋值为\co{NULL}时或该指针尚未被读者访问时
替代\co{rcu_assign_pointer()}使用。
使用\co{RCU_INIT_POINTER()}为编译器在选择指令和
执行优化方面提供了更多余地，从而提高了性能。

另一方面，错误使用\co{RCU_INIT_POINTER()}可能导致
无声的内存损坏，所以请务必小心！
是的，在某些情况下，内核可以检查在给定内核上下文中
是否不当使用了RCU API成员，但
\co{RCU_INIT_POINTER()}使用的约束尚无法被检查。

最后，在Linux内核中，上述人类认知的局限性
因Linux上运行的工作负载的多样性和严苛性而加剧。
截至v5.16，这已产生了不少于五种RCU风格，
每种都旨在为RCU读者和写者提供不同的性能、
可扩展性、响应时间和能效权衡。
这些RCU风格是下一节的主题。

\subsubsection{RCU拥有一系列等待完成API}
\label{sec:defer:RCU has a Family of Wait-to-Finish APIs}

对``什么是RCU''最直接的回答是RCU是一个API。
例如，Linux内核中使用的RCU实现由
\cref{tab:defer:RCU Wait-to-Finish APIs}概括，
其各列分别显示了RCU、"可睡眠"RCU（SRCU）、Tasks RCU、
Tasks RCU Rude、Tasks RCU Trace和通用API的等待读者部分，
\cref{tab:defer:RCU Markers for Read-Side Critical Sections}
显示了RCU读侧标记的完整范围，
\cref{tab:defer:SRCU Markers for Read-Side Critical Sections}
显示了SRCU读侧标记的完整范围，
\cref{tab:defer:RCU Polled Update-Side Primitives}
显示了完整的RCU轮询grace period API，
\cref{tab:defer:RCU Publish-Subscribe and Version Maintenance APIs}
显示了API~\cite{PaulEMcKenney2024RCUAPI}的发布-订阅部分。\footnote{
	此引用涵盖v6.10及更高版本。
	早期版本的Linux内核RCU API文档可在其他地方找到~\cite{PaulEMcKenney2008WhatIsRCUAPI,PaulEMcKenney2014RCUAPI,PaulEMcKenney2019RCUAPI}。}

\begin{sidewaystable*}[tbp]
\rowcolors{1}{}{lightgray}
\renewcommand*{\arraystretch}{1.3}
\centering
\caption{RCU等待完成API}
\label{tab:defer:RCU Wait-to-Finish APIs}
\scriptsize\IfEbookSize{\hspace*{-1.8in}}{\hspace*{-.125in}}
\ebresizeverb{0.7}{
\begin{tabularx}{8.5in}{>{\raggedright\arraybackslash}p{0.92in}
    >{\raggedright\arraybackslash}p{1.37in}
    >{\raggedright\arraybackslash}p{1.39in}
    >{\raggedright\arraybackslash}p{1.09in}
    >{\raggedright\arraybackslash}p{1.34in}
    >{\raggedright\arraybackslash}p{1.44in}}
\toprule
&
    {\bf RCU}：原始版 &
	{\bf SRCU}：可睡眠读者 &
	    {\bf Tasks RCU}：释放追踪蹦床 &
		{\bf Tasks RCU Rude}：释放空闲任务追踪蹦床 &
		    {\bf Tasks RCU Trace}：保护可睡眠BPF程序 \\
\midrule
{\bf 初始化和清理} &
    &
	\tco{DEFINE_SRCU()} \tco{DEFINE_STATIC_SRCU()}
	\tco{init_srcu_struct()} \tco{cleanup_srcu_struct()} &
	    &
		&
		    \\
{\bf Read-side critical section标记} &
    \tco{rcu_read_lock()}~! \tco{rcu_read_unlock()}~!  ~~~~~~~~~~~~~~~~~~
    （更多参见\cref{tab:defer:RCU Markers for Read-Side Critical Sections}。） &
	\tco{srcu_read_lock()} \tco{srcu_read_unlock()} ~~~~~~~~~~~~~~~~~
	（更多参见\cref{tab:defer:SRCU Markers for Read-Side Critical Sections}。） &
	    非空闲任务中的自愿上下文切换。 &
		禁用抢占的代码区域。 &
		    \tco{rcu_read_lock_trace()} \tco{rcu_read_unlock_trace()}
		    \tco{guard(rcu_tasks_trace)()}
		    \tco{scoped_guard(rcu_tasks_trace)} \\
{\bf 更新侧原语（同步） } &
    { \tco{synchronize_rcu()}
      \tco{synchronize_net()}
      \tco{synchronize_rcu_expedited()} } &
	\tco{synchronize_srcu()} \tco{synchronize_srcu_expedited()} &
	    \tco{synchronize_rcu_tasks()}
	    \tco{synchronize_rcu_mult()} &
		\tco{synchronize_rcu_tasks_rude()} &
		    \tco{synchronize_rcu_tasks_trace()}
		    \tco{rcu_trace_implies_rcu_gp()} \\
{\bf 更新侧原语（异步/callback） } &
    \tco{call_rcu()} ! ~~~~~~
    \tco{call_rcu_hurry()} &
	\tco{call_srcu()} &
	    \tco{call_rcu_tasks()} &
		\tco{call_rcu_tasks_rude()} &
		    \tco{call_rcu_tasks_trace()} \\
{\bf 更新侧原语（等待callback） } &
    \tco{rcu_barrier()} &
	\tco{srcu_barrier()} &
	    \tco{rcu_barrier_tasks()} &
		\tco{rcu_barrier_tasks_rude()} &
		    \tco{rcu_barrier_tasks_trace()} \\
{\bf 更新侧原语（轮询） } &
    \tco{get_state_synchronize_rcu()}
    \tco{start_poll_synchronize_rcu()}
    \tco{poll_state_synchronize_rcu()} ~~~~~~~~~~~~~~~~~~
    （更多参见\cref{tab:defer:RCU Polled Update-Side Primitives}。） &
	\tco{get_state_synchronize_srcu()}
	\tco{start_poll_synchronize_srcu()}
	\tco{poll_state_synchronize_srcu()} &
	    &
		&
		    \\
{\bf 更新侧原语（发起/等待）} &
    \tco{get_state_synchronize_rcu()}
    \tco{cond_synchronize_rcu()} &
	&
	    &
		&
		    \\
{\bf 更新侧原语（释放内存） } &
    \tco{kfree_rcu()}
    \tco{kfree_rcu_mightsleep()}
    \tco{kvfree_rcu()}
    \tco{kvfree_rcu_mightsleep()} &
	&
	    &
		&
		    \\
{\bf 类型安全内存 } &
    \tco{SLAB_TYPESAFE_BY_RCU} &
	&
	    &
		&
		    \\
{\bf 读侧约束 } &
    不可阻塞（仅可抢占） &
	不可对同一\tco{srcu_struct}使用\tco{synchronize_srcu()} &
	    不可自愿上下文切换 &
		既不可阻塞也不可抢占 &
			不可使用RCU tasks trace grace period \\
{\bf 读侧开销 } &
    CPU本地访问（\tco{PREEMPT=n}时为\tco{barrier()}） &
	简单指令、内存屏障 &
	    免费 &
		CPU本地访问（\tco{PREEMPT=n}时免费） &
		    CPU本地访问 \\
{\bf 异步更新侧开销 } &
    亚微秒级 &
	亚微秒级 &
	    亚微秒级 &
		亚微秒级 &
		    亚微秒级 \\
{\bf Grace period延迟 } &
    数十毫秒 &
        毫秒级 &
	    秒级 &
		毫秒级 &
		    数十毫秒 \\
{\bf 加速grace period延迟 } &
    数十微秒 &
        微秒级 &
	    不适用 &
		不适用 &
		    不适用 \\
\bottomrule
\end{tabularx}
}
\end{sidewaystable*}

\begin{table*}
\rowcolors{1}{}{lightgray}
\renewcommand*{\arraystretch}{1.15}
\footnotesize
\centering\OneColumnHSpace{-.4in}
\ebresizewidth{
\begin{tabular}{llp{2.0in}}
\toprule
获取 &
	释放 &
		说明 \\
\midrule
\tco{rcu_read_lock()} &
	\tco{rcu_read_unlock()} & \\
\tco{guard(rcu)()} & & RCU读者持续到外围作用域结束。 \\
\tco{scoped_guard(rcu)} & & RCU读者用于紧随其后的语句。 \\
\tco{rcu_read_lock_bh()} &
	\tco{rcu_read_unlock_bh()} & \\
\tco{rcu_read_lock_sched()} &
	\tco{rcu_read_unlock_sched()} & \\
\tco{rcu_read_lock_sched_notrace()} &
	\tco{rcu_read_unlock_sched_notrace()} & \\
\tco{local_bh_disable()} &
	\tco{local_bh_enable()} &
		以及其他任何禁用下半部的操作。 \\
\tco{preempt_disable()} &
	\tco{preempt_enable()} &
		以及其他任何禁用抢占的操作。 \\
\tco{local_irq_save()} &
	\tco{local_irq_restore()} &
		以及其他任何禁用中断的操作。 \\
\bottomrule
\end{tabular}
}
\caption{RCU Read-Side Critical Section标记}
\label{tab:defer:RCU Markers for Read-Side Critical Sections}
\end{table*}

\begin{table*}
\rowcolors{1}{}{lightgray}
\renewcommand*{\arraystretch}{1.15}
\footnotesize
\centering\OneColumnHSpace{-.4in}
\ebresizewidth{
\begin{tabular}{llp{2.0in}}
\toprule
获取 &
	释放 &
		说明 \\
\midrule
\tco{srcu_read_lock()} &
	\tco{srcu_read_unlock()} & \\
\tco{guard(srcu)(&my_srcu)} & & SRCU读者持续到外围作用域结束。 \\
\tco{scoped_guard(srcu, &my_srcu)} & & SRCU读者用于紧随其后的语句。 \\
\tco{srcu_read_lock_nmisafe()} &
	\tco{srcu_read_unlock_nmisafe()} & \\
\tco{srcu_read_lock_notrace()} &
	\tco{srcu_read_unlock_notrace()} & \\
\tco{srcu_down_read()} &
	\tco{srcu_up_read()} & \\
\bottomrule
\end{tabular}
}
\caption{SRCU Read-Side Critical Section标记}
\label{tab:defer:SRCU Markers for Read-Side Critical Sections}
\end{table*}

\begin{table*}
\rowcolors{1}{}{lightgray}
\renewcommand*{\arraystretch}{1.15}
\footnotesize
\centering\OneColumnHSpace{-.4in}
\ebresizewidth{
\begin{tabular}{p{2.5in}p{2.5in}}
\toprule
压缩状态 / 完整状态 &
		说明 \\
\midrule
\tco{get_completed_synchronize_rcu()}
\tco{get_completed_synchronize_rcu_full()} &
	返回始终已完成的状态。 \\
\tco{get_state_synchronize_rcu()}
\tco{get_state_synchronize_rcu_full()} &
	返回用于grace period完成检测的状态。 \\
\tco{start_poll_synchronize_rcu()}
\tco{start_poll_synchronize_rcu_full()} &
	返回用于grace period完成检测的状态，
	并在需要时启动grace period。 \\
\tco{start_poll_synchronize_rcu_expedited()}
\tco{start_poll_synchronize_rcu_expedited_full()} &
	返回用于grace period完成检测的状态，
	并在需要时启动加速的grace period。 \\
\tco{poll_state_synchronize_rcu()}
\tco{poll_state_synchronize_rcu_full()} &
	如果该状态的grace period已结束则返回true。 \\
\tco{cond_synchronize_rcu()}
\tco{cond_synchronize_rcu_full()} &
	（如需要）等待该状态的grace period结束。 \\
\tco{cond_synchronize_rcu_expedited()}
\tco{cond_synchronize_rcu_expedited_full()} &
	（如需要）以加速方式等待该状态的grace period结束。 \\
\tco{NUM_ACTIVE_RCU_POLL_OLDSTATE}
\tco{NUM_ACTIVE_RCU_POLL_OLDSTATE} &
	对应进行中grace period的最大状态数。 \\
\tco{same_state_synchronize_rcu()}
\tco{same_state_synchronize_rcu_full()} &
	两个状态是否等价？ \\
\bottomrule
\end{tabular}
}
\caption{RCU轮询更新侧原语}
\label{tab:defer:RCU Polled Update-Side Primitives}
\end{table*}

如果你是RCU新手，可以考虑只关注
\cref{tab:defer:RCU Wait-to-Finish APIs}
中的某一列，每一列概述了Linux内核RCU API系列的一个成员。
例如，如果你主要对理解RCU在Linux内核中的使用感兴趣，
``RCU''列将是起点，因为它使用得最频繁。
另一方面，如果你想为了理解RCU本身而学习，
``Tasks RCU Rude''拥有最简单的API。
你随时可以回来看其他列。

如果你已经熟悉RCU，这些表格可以
作为有用的参考。

\QuickQuiz{
	为什么\cref{tab:defer:RCU Wait-to-Finish APIs}
	中的一些单元格有感叹号（``!''）？
}\QuickQuizAnswer{
	带有感叹号的API成员（\co{rcu_read_lock()}、
	\co{rcu_read_unlock()}和\co{call_rcu()}）是
	Paul E. McKenney在90年代中期所知的
	Linux RCU API的仅有成员。
	在那个时期，他错误地认为自己
	已经了解了关于RCU的一切。
}\QuickQuizEnd

``RCU''列对应于三个Linux内核RCU
实现的合并~\cite{PaulEMcKenney2019RCUCVE,McKenney:2019:CRS:3319647.3325836}，
其中RCU read-side critical section以
\apik{rcu_read_lock()}开始、以\apik{rcu_read_unlock()}结束。
这些临界区也可以使用
\cref{tab:defer:RCU Markers for Read-Side Critical Sections}中的API，包括
\apik{rcu_read_lock_bh()}、\apik{rcu_read_lock_sched()}、
\apik{rcu_read_unlock_bh()}或\apik{rcu_read_unlock_sched()}。
此外，任何禁用下半部、中断或抢占的代码区域
也充当RCU read-side critical section。
RCU read-side critical section可以嵌套。\footnote{
	请注意\co{guard(rcu)()}和\co{scoped_guard(rcu)}
	提供的RAII变体。}
对应的同步更新侧原语
\apik{synchronize_rcu()}和\apik{synchronize_rcu_expedited()}，
以及它们的同义词\apik{synchronize_net()}，
等待任何类型的当前正在执行的RCU read-side critical section完成。
这种等待的时长被称为``\IX{grace period}''，
\apik{synchronize_rcu_expedited()}旨在以增加CPU开销和IPI
为代价减少\IXh{grace-period}{latency}。
异步更新侧原语\apik{call_rcu()}在
后续grace period之后调用指定的函数和指定的参数。
例如，\co{call_rcu(p,f);}将导致
``RCU callback'' \co{f(p)}在后续grace period之后被调用。
使用\co{CONFIG_RCU_LAZY=y}构建的内核会将
\co{call_rcu()}本来会立即启动的grace period延迟最多数秒。
在这种延迟不可接受的情况下，可以使用
\apik{call_rcu_hurry()}来立即启动grace period。
有些情况下，例如卸载使用\co{call_rcu()}或
\co{call_rcu_hurry()}的Linux内核模块时，
需要等待所有未完成的RCU callback完成~\cite{PaulEMcKenney2007rcubarrier}。
\apik{rcu_barrier()}原语完成此任务。

\QuickQuizSeries{%
\QuickQuizB{
	如何防止大量RCU read-side critical section
	无限期地阻塞\co{synchronize_rcu()}调用？
}\QuickQuizAnswerB{
	无需做任何事情来防止RCU read-side critical section
	无限期地阻塞\co{synchronize_rcu()}调用，
	因为\co{synchronize_rcu()}调用只需等待
	\emph{已有的}RCU read-side critical section。
	因此，只要每个RCU read-side critical section
	的持续时间是有限的，RCU grace period也将保持有限。
}\QuickQuizEndB
%
\QuickQuizM{
	\co{synchronize_rcu()} API会等待所有已有的
	中断处理程序完成，对吧？
}\QuickQuizAnswerM{
	在v4.20及更高版本的Linux内核中，
	是的~\cite{PaulEMcKenney2019RCUCVE,McKenney:2019:CRS:3319647.3325836}。

	但在更早的内核中不是这样，尤其是在使用
	可抢占RCU时！
	你应当改用\apik{synchronize_irq()}。
	或者，你可以在希望\co{synchronize_rcu()}等待的
	特定中断处理程序中放置\co{rcu_read_lock()}和
	\co{rcu_read_unlock()}调用。
	但即使如此也要小心，因为可抢占RCU不保证
	等待中断处理程序中\co{rcu_read_lock()}之前或
	\co{rcu_read_unlock()}之后的部分。
}\QuickQuizEndM
%
\QuickQuizM{
	既然任何禁用中断的代码区域都充当
	RCU read-side critical section，那么原子操作（atomic operation）和
	单条机器指令是否也充当微小的RCU read-side
	critical section？
}\QuickQuizAnswerM{
	是的，确实如此。

	但如果你想使用这个技巧，请三思，
	而且这次要更认真地思考。

	如果你确实别无选择只能使用此技巧，
	请仔细记录你对它的使用。
	否则，阅读代码的每个人（包括你自己）都会
	因此而恨你。

	为什么如此多的仇恨？

	要理解原因，请记住Jack Slingwine和我发明RCU的
	一大优势在于我们标记了RCU读者。
	听听所有那些独立发明了类似RCU东西但
	未能标记其read-side critical section的人的话：
	你确实需要标记你的RCU读者！
	包括你的单指令RCU读者。
}\QuickQuizEndM
%
\QuickQuizE{
	\co{synchronize_rcu()}和\co{rcu_barrier()}
	之间有什么区别？
}\QuickQuizAnswerE{
	它们等待的东西不同。
	\co{synchronize_rcu()}等待已有的RCU read-side
	critical section完成，而\co{rcu_barrier()}等待
	先前\co{call_rcu()}调用的callback被调用。

\begin{table}
\renewcommand*{\arraystretch}{1.2}
\centering
\small
\tcresizewidth{
\begin{fcvref}[ln:defer:synchonize-rcu vs rcu-barrier]
\begin{tabular}{lll}
\toprule
            & \multicolumn{2}{c}{必须等待直到（代码片段中的）：} \\
\cmidrule{2-3}
\multicolumn{1}{c}{调用于：} & \multicolumn{1}{c}{\tco{synchronize_rcu()}}
					& \multicolumn{1}{c}{\tco{rcu_barrier()}} \\
\cmidrule(r){1-1} \cmidrule{2-3}
\tco{do_something_1()} & 			 & \\
\tco{do_something_2()} & \tco{rcu_read_unlock()} (\clnref{rrul}) & \\
\tco{do_something_3()} & \tco{rcu_read_unlock()} (\clnref{rrul})
						 & \tco{f(&p->rh)} (\clnref{cb}) \\
\tco{do_something_4()} &			 & \tco{f(&p->rh)} (\clnref{cb}) \\
\tco{do_something_5()} & 			 & \\
\bottomrule
\end{tabular}
\end{fcvref}
}

\begin{fcvlabel}[ln:defer:synchonize-rcu vs rcu-barrier]
\begin{VerbatimN}[commandchars=\\\@\$]
	do_something_1();			\lnlbl@ds1$
	rcu_read_lock();			\lnlbl@rrl$
	do_something_2();			\lnlbl@ds2$
	call_rcu(&p->rh, f);			\lnlbl@cr$
	do_something_3();			\lnlbl@ds3$
	rcu_read_unlock();			\lnlbl@rrul$
	do_something_4();			\lnlbl@ds4$
	// f(&p->rh) invoked			\lnlbl@cb$
	do_something_5();			\lnlbl@ds5$
\end{VerbatimN}
\end{fcvlabel}
\caption{\tco{synchronize_rcu()} vs. \tco{rcu_barrier()}}
\label{tab:defer:synchonize-rcu vs rcu-barrier}
\end{table}

	\cref{tab:defer:synchonize-rcu vs rcu-barrier}
	说明了这一区别，其中底部的代码片段展示了
	给定CPU正在执行的代码。
	为简单起见，假设没有其他CPU在执行
	\co{rcu_read_lock()}、\co{rcu_read_unlock()}或
	\co{call_rcu()}。

	该表展示了如果与各\co{do_something_*()}
	函数并发调用，每个原语必须等待多久，
	空单元格表示无需等待。
	如你所见，\co{synchronize_rcu()}只有在
	RCU read-side critical section中才需要等待，
	在这种情况下它必须等待结束该临界区的
	\co{rcu_read_unlock()}。
	相反，RCU read-side critical section对
	\co{rcu_barrier()}没有影响。
	然而，当\co{rcu_barrier()}在\co{call_rcu()}调用之后
	执行时，它必须等到相应的RCU callback被调用。

	话虽如此，有一种特殊情况，每次
	\co{rcu_barrier()}调用都可以被直接的
	\co{synchronize_rcu()}调用替换，即
	\co{synchronize_rcu()}是用\co{call_rcu()}实现的，
	且只有一个全局callback列表的情况。
	但请勿在可移植代码中这样做！！！
}\QuickQuizEndE
}

\co{synchronize_rcu()}和\co{call_rcu()}系列API易于使用，
自动为用户提供排队、批处理和唤醒功能。
然而，有些情况下这种易用性反而成为障碍，
例如，当受RCU保护的对象缓存可能在本应允许
释放它们的grace period期间被重用时。
虽然这可以通过仔细使用与\co{call_rcu()}排队的
RCU callback交互的标志来处理，但这可能不便，
并且由于注定失败的\co{call_rcu()}调用的开销而浪费CPU时间。

在这些情况下，RCU的轮询grace period原语可以提供帮助。
\apik{get_state_synchronize_rcu()}和\apik{start_poll_synchronize_rcu()}
函数返回一个grace period状态"cookie"，它封装了
后续grace period完成后的grace period状态。
一旦这样的后续grace period结束，
\apik{poll_state_synchronize_rcu()}将返回\co{true}。
RCU的轮询grace period API中还有相当多的其他函数。
有关它们的更多信息，包括示例代码，请参见
2024年LSF/MM关于此主题的演示幻灯片第34页及
之后的内容~\cite{PaulEMcKenney2024ReclamationInteractionsWithRCU}。\footnote{
	Linux内核源代码中这些函数的kernel-doc头
	也可能非常有帮助。}

最后，RCU可用于提供
\IX{type-safe memory}~\cite{Cheriton96a}，
如\cref{sec:defer:Type-Safe Memory}所述。
在RCU的上下文中，类型安全内存保证给定的
数据元素在访问它的任何RCU read-side critical section
期间不会改变类型。
要使用基于RCU的类型安全内存，请将
\apik{SLAB_TYPESAFE_BY_RCU}传递给\apik{kmem_cache_create()}。

\cref{tab:defer:RCU Wait-to-Finish APIs}中的``SRCU''列
展示了一个专用的RCU API，允许在SRCU read-side
critical section~\cite{PaulEMcKenney2006c}中进行通用的睡眠，
这些临界区由\apik{srcu_read_lock()}和\apik{srcu_read_unlock()}界定。
然而，与RCU不同，SRCU的\apik{srcu_read_lock()}返回一个值，
必须将其传入对应的\apik{srcu_read_unlock()}。
这一区别是因为SRCU用户为每种不同的SRCU用法分配一个
\apik{srcu_struct}，因此没有方便的地方存储每任务的
读者嵌套计数。
（请记住，虽然Linux内核提供动态分配的
每CPU存储，但目前还没有动态分配的每任务存储。）
SRCU read-side critical section也可以用
\cref{tab:defer:SRCU Markers for Read-Side Critical Sections}中显示的API标记，
同样注意新的RAII变体。

给定的\co{srcu_struct}结构可以用\co{DEFINE_SRCU()}定义为
全局变量（如果该结构必须在多个翻译单元中使用），
或用\co{DEFINE_STATIC_SRCU()}定义。
例如，\co{DEFINE_SRCU(my_srcu)}将创建一个名为
\co{my_srcu}的全局变量，可被程序中的任何文件使用。
或者，\co{srcu_struct}结构可以是栈上变量
或动态分配的内存区域。
在这两种非全局变量的情况下，内存必须在首次使用前
用\co{init_srcu_struct()}初始化，
并在最后一次使用后（但在底层存储消失之前）
用\co{cleanup_srcu_struct()}清理。

无论如何创建，这些不同的\co{srcu_struct}结构
防止SRCU read-side critical section阻塞
不相关的\apik{synchronize_srcu()}和
\apik{synchronize_srcu_expedited()}调用。
当然，在SRCU read-side critical section内使用
\apik{synchronize_srcu()}或\apik{synchronize_srcu_expedited()}
可能导致自死锁，因此应当避免（v6.4及更高版本中
lockdep会检查这一点）。
与RCU一样，SRCU的\co{synchronize_srcu_expedited()}比
\co{synchronize_srcu()}减少了grace period延迟，
但代价是增加CPU开销。
同样与RCU类似，有\apik{get_state_synchronize_srcu()}、
\apik{start_poll_synchronize_srcu()}和
\apik{poll_state_synchronize_srcu()}函数允许轮询SRCU grace period API。

\QuickQuiz{
	在什么条件下可以在SRCU read-side critical section内
	安全地使用\co{synchronize_srcu()}？
}\QuickQuizAnswer{
	原则上，你可以在使用某个\co{srcu_struct}的
	SRCU read-side critical section内，
	对另一个\co{srcu_struct}使用\co{synchronize_srcu()}或
	\co{synchronize_srcu_expedited()}。
	然而，在实践中这样做几乎肯定是个坏主意。
	特别是，\cref{lst:defer:Multistage SRCU Deadlocks}
	中所示的代码仍然可能导致死锁，
	在v6.4及更高版本的Linux内核中lockdep会对此发出警告。

\begin{listing}
\begin{VerbatimL}
idx = srcu_read_lock(&ssa);
synchronize_srcu(&ssb);
srcu_read_unlock(&ssa, idx);

/* . . . */

idx = srcu_read_lock(&ssb);
synchronize_srcu(&ssa);
srcu_read_unlock(&ssb, idx);
\end{VerbatimL}
\caption{多阶段SRCU死锁}
\label{lst:defer:Multistage SRCU Deadlocks}
\end{listing}
}\QuickQuizEnd

与普通RCU类似，可以使用异步的\apik{call_srcu()}函数来
避免自死锁。
然而，使用\co{call_srcu()}时必须特别小心，
因为单个任务可能非常快速地注册SRCU callback。
鉴于SRCU读者可以阻塞任意长的时间，
这可能消耗任意大量的内存。
相比之下，使用同步的\co{synchronize_srcu()}
接口，给定任务必须完成等待当前grace period，
然后才能开始等待下一个。

同样与RCU类似，有一个\apik{srcu_barrier()}函数，
等待所有先前\co{call_srcu()} callback被调用。

换言之，SRCU通过允许开发者限制其作用域来
弥补其极弱的前向进展保证。

\cref{tab:defer:RCU Wait-to-Finish APIs}中的``Tasks RCU''列
展示了一个专用RCU API，用于协调Linux内核追踪中
使用的蹦床的释放。
这些蹦床用于将控制从被追踪代码中的某一点
转移到执行实际追踪的代码。
当然，必须确保在释放给定蹦床之前，
该蹦床中所有正在执行的代码都已完成。

对被追踪代码的更改通常仅限于单条跳转或调用指令，
因此无法容纳实现\co{rcu_read_lock()}和
\co{rcu_read_unlock()}所需的代码序列。
蹦床本身也不能包含这些\co{rcu_read_lock()}和
\co{rcu_read_unlock()}调用。
要理解这一点，考虑一个即将开始执行
给定蹦床的CPU。
因为它尚未执行\co{rcu_read_lock()}，该
蹦床可能随时被释放，这对该CPU来说将是致命的
意外。
因此，蹦床不能通过在被追踪代码或蹦床本身中
执行的同步原语来保护。
这就引出了蹦床究竟如何被保护的问题。

回答这个问题的关键在于注意到蹦床代码
从不包含直接或间接进行自愿上下文切换的代码。
此代码可能被抢占，但它永远不会直接或间接
调用\apik{schedule()}。
这提示了一种RCU变体，其中自愿上下文切换和
空闲执行是唯一的quiescent state。
此变体就是Tasks RCU。

Tasks RCU的不寻常之处在于没有读侧标记函数，
鉴于其主要用例无处放置此类标记，这倒是件好事。
相反，\co{schedule()}的调用直接充当quiescent state，
但仅限于那些对应自愿上下文切换的\co{schedule()}调用。
更新可以使用\apik{synchronize_rcu_tasks()}来等待所有
已有的蹦床执行完成，或者使用其异步对应物
\apik{call_rcu_tasks()}。
还有一个\apik{rcu_barrier_tasks()}，等待所有先前
\co{call_rcu_tasks()}调用对应的callback完成。
没有\co{synchronize_rcu_tasks_expedited()}，因为
尚未有此需求，不过实现一个有用的变体也并非
没有挑战。

\QuickQuiz{
	在使用\co{CONFIG_PREEMPT_NONE=y}构建的内核中，
	\co{synchronize_rcu()}不是会等待所有蹦床完成吗？
	毕竟抢占被禁用了，而蹦床从不直接或间接
	调用\co{schedule()}。
}\QuickQuizAnswer{
	你说得对！

	实际上，在不可抢占内核中，\co{synchronize_rcu_tasks()}
	是\co{synchronize_rcu()}的包装器。
	但请注意，具有运行时可设置抢占
	（\co{CONFIG_PREEMPT_DYNAMIC=y}）或启用了惰性抢占
	（\co{CONFIG_PREEMPT_LAZY=y}）的内核实际上是可抢占的。
}\QuickQuizEnd

``Tasks RCU Rude''列提供了
\cref{sec:defer:Toy Implementation}中展示的
玩具实现的更有效变体。
此变体使每个CPU执行一次上下文切换，
这样任何自愿上下文切换或任何可抢占的代码区域
都可以充当quiescent state。
Tasks RCU Rude变体使用Linux内核的workqueue设施来
强制并发上下文切换，与玩具实现采取的
串行逐CPU方法形成对比。
其API与Tasks RCU类似，包括缺少显式的读侧标记。

最后，``Tasks RCU Trace''列提供了一个
功能类似于SRCU的RCU实现，但具有快得多的
读侧标记。\footnote{
	因此在Tasks RCU家族中是不寻常的，因为它有
	显式的读侧标记！}
然而，这种速度是因为这些标记不执行内存屏障指令，
这意味着Tasks RCU Trace grace period通常必须向
所有CPU发送IPI，从而降低实时响应。
尽管如此，在没有读者的情况下，由此产生的
grace period延迟相当短，可与RCU相媲美。

\subsubsection{RCU拥有发布-订阅和版本维护API}
\label{sec:defer:RCU has Publish-Subscribe and Version-Maintenance APIs}

幸运的是，\cref{tab:defer:RCU Publish-Subscribe and Version Maintenance APIs}
中所示的RCU发布-订阅和版本维护原语
适用于上面讨论的所有RCU变体。
这种共性允许更多代码被共享，并减少API的增殖。
RCU发布-订阅API的最初目的是将内存屏障
嵌入这些API中，使Linux内核程序员可以使用RCU
而无需成为Linux支持的20多种CPU系列中每一种的
内存排序模型的专家~\cite{Spraul01}。

\begin{table*}
\renewcommand*{\arraystretch}{1.15}
\footnotesize
\centering\OneColumnHSpace{-.4in}
\ebresizewidth{
\begin{tabular}{llp{2.2in}}
\toprule
类别 &
	原语 &
		开销 \\
\midrule
指针发布 &
	\tco{rcu_assign_pointer()} &
		内存屏障 \\
&
	\tco{rcu_replace_pointer()} &
		内存屏障（Alpha上两个） \\
&
	\tco{rcu_pointer_handoff()} &
		简单指令 \\
&
	\tco{unrcu_pointer()} &
		简单指令 \\
&
	\tco{RCU_INIT_POINTER()} &
		简单指令 \\
&
	\tco{RCU_POINTER_INITIALIZER()} &
		编译时常量 \\
\midrule
指针订阅（遍历） &
	\tco{rcu_access_pointer()} &
		简单指令 \\
&
	\tco{rcu_dereference()} &
		简单指令（Alpha上为内存屏障） \\
&
	\tco{rcu_dereference_check()} &
		简单指令（Alpha上为内存屏障） \\
&
	\tco{rcu_dereference_bh()} &
		简单指令（Alpha上为内存屏障） \\
&
	\tco{rcu_dereference_bh_check()} &
		简单指令（Alpha上为内存屏障） \\
&
	\tco{rcu_dereference_sched()} &
		简单指令（Alpha上为内存屏障） \\
&
	\tco{rcu_dereference_sched_check()} &
		简单指令（Alpha上为内存屏障） \\
&
	\tco{rcu_dereference_protected()} &
		简单指令 \\
&
	\tco{rcu_dereference_raw()} &
		简单指令（Alpha上为内存屏障） \\
&
	\tco{rcu_dereference_raw_check()} &
		简单指令（Alpha上为内存屏障） \\
\midrule
SRCU指针订阅（遍历） &
	\tco{srcu_dereference()} &
		简单指令 \\
&
	\tco{srcu_dereference_check()} &
		简单指令 \\
&
	\tco{srcu_dereference_notrace()} &
		简单指令 \\
\bottomrule
\end{tabular}
}
\caption{RCU发布-订阅和版本维护API}
\label{tab:defer:RCU Publish-Subscribe and Version Maintenance APIs}
\end{table*}

这些原语直接操作指针，对于创建受RCU保护的
链式数据结构（如受RCU保护的数组和树）非常有用。
链表的特殊情况由
\cref{sec:defer:RCU has List-Processing APIs}
中描述的一组单独的API处理。

第一类发布指向新数据项的指针。
\apik{rcu_assign_pointer()}原语确保在弱序机器上，
任何先前的初始化保持排在指针赋值之前。
\apik{rcu_replace_pointer()}原语像\co{rcu_assign_pointer()}
一样更新指针，但也像\co{rcu_dereference_protected()}
（见下文）那样返回先前的值，包括lockdep表达式。
当更新者必须同时发布新指针并释放旧指针引用的结构时，
这种替换非常方便。

\QuickQuizSeries{%
\QuickQuizB{
	通常，任何受\co{rcu_dereference()}保护的指针
	\emph{必须}始终使用
	\cref{tab:defer:RCU Publish-Subscribe and Version Maintenance APIs}
	中的指针发布函数来更新，
	例如\co{rcu_assign_pointer()}。

	有没有例外？
}\QuickQuizAnswerB{
	一个例外是：当多元素链式数据结构
	在其他CPU不可访问时作为一个整体初始化，
	然后用单个\co{rcu_assign_pointer()}
	设置一个指向此数据结构的全局指针。
	初始化时的指针赋值不需要使用
	\co{rcu_assign_pointer()}，
	尽管在结构全局可见后发生的任何赋值
	\emph{必须}使用\co{rcu_assign_pointer()}。

	然而，除非此初始化代码处于非常热的代码路径上，
	即使理论上不必要，使用\co{rcu_assign_pointer()}
	可能也是明智的。
	一个"小"变更使你珍视的关于初始化
	私下进行的假设失效太容易了。
}\QuickQuizEndB
%
\QuickQuizE{
	这些遍历和更新原语可以与任何RCU API系列成员
	一起使用，这有什么缺点吗？
}\QuickQuizAnswerE{
	有时自动化代码检查器如``sparse''
	（或者实际上人类也是如此）
	很难确定给定的RCU遍历原语对应于
	哪种类型的RCU read-side critical section。
	例如，考虑\cref{lst:defer:Diverse RCU Read-Side Nesting}
	中所示的代码。

\begin{listing}
\begin{VerbatimL}
rcu_read_lock();
preempt_disable();
p = rcu_dereference(global_pointer);

/* . . . */

preempt_enable();
rcu_read_unlock();
\end{VerbatimL}
\caption{多样化的RCU读侧嵌套}
\label{lst:defer:Diverse RCU Read-Side Nesting}
\end{listing}

	\co{rcu_dereference()}原语是在普通的RCU
	临界区中还是在RCU Sched临界区中？
	你需要做什么才能弄清楚？

	但也许在v4.20 Linux内核中合并RCU风格之后，
	我们不再需要关心了！
}\QuickQuizEndE
}

\apik{rcu_pointer_handoff()}原语简单地返回其唯一参数，
但对于检查从RCU read-side critical section泄漏指针的
工具很有用。
使用\co{rcu_pointer_handoff()}向此类工具表明，
所讨论结构的保护已从RCU移交给某种其他机制，
如加锁或引用计数。

\apik{RCU_INIT_POINTER()}宏可用于初始化
尚未暴露给读者的受RCU保护的指针，
或者将受RCU保护的指针设为\co{NULL}。
在这些受限的情况下，不需要\co{rcu_assign_pointer()}
提供的内存屏障指令。
类似地，\apik{RCU_POINTER_INITIALIZER()}提供了一个
\GCC 风格的结构初始化器，以便于在结构中
初始化受RCU保护的指针。

第二类订阅指向数据项的指针，或者说
安全地遍历受RCU保护的指针。
同样，简单地使用C语言访问加载这些指针
可能导致看到被指向数据中的初始化前垃圾。
类似地，在RCU read-side critical section之外
以任何方式加载这些指针都可能导致被指向的对象
随时被释放。
然而，如果指针仅用于测试而不进行解引用，
被指向对象的释放不一定是问题。
在这种情况下，可以使用\apik{rcu_access_pointer()}。
但通常需要RCU读侧保护，因此
\apik{rcu_dereference()}原语使用Linux内核的\co{lockdep}
设施~\cite{JonathanCorbet2006lockdep}来验证此
\co{rcu_dereference()}调用受到
\co{rcu_read_lock()}、\co{srcu_read_lock()}或其他RCU读侧
标记的保护。
相比之下，\co{rcu_access_pointer()}原语不涉及
\co{lockdep}，因此在RCU read-side critical section
之外使用时不会引发\co{lockdep}警告。

另一种不需要保护的情况是更新侧代码
在持有更新侧锁的同时访问受RCU保护的指针。
\apik{rcu_dereference_protected()} API成员就是为此情况
提供的。其第一个参数是受RCU保护的指针，
第二个参数接受一个lockdep表达式，
描述必须持有哪些锁才能使访问安全。
从读者和更新者都调用的代码可以使用
\apik{rcu_dereference_check()}，它也接受lockdep表达式，
但也可以从不持有锁的读侧代码中调用。
在某些情况下，lockdep表达式可能非常复杂，
例如，使用细粒度加锁时，
大量锁中的任何一个都可能被持有，
要确定哪个适用可能相当困难。
在这些（希望罕见的）情况下，\apik{rcu_dereference_raw()}
提供保护但不检查是否在读者内调用或持有任何特定的锁。
\apik{rcu_dereference_raw_notrace()} API成员行为类似，
但不可被追踪，因此可以安全地被追踪代码使用。

虽然几乎任何链式结构都可以通过操作指针来访问，
但更高层的结构可以非常有帮助。
因此下一节查看Linux内核使用的各种
受RCU保护的链表。

\subsubsection{RCU拥有链表处理API}
\label{sec:defer:RCU has List-Processing APIs}

\begin{figure}
\centering
\resizebox{3in}{!}{\includegraphics{defer/Linux_list}}
\caption{Linux环形双向链表（\tco{list}）}
\label{fig:defer:Linux Circular Linked List (list)}
\end{figure}

\begin{figure}
\centering
\resizebox{3in}{!}{\includegraphics{defer/Linux_list_abbr}}
\caption{Linux链表简略图}
\label{fig:defer:Linux Linked List Abbreviated}
\end{figure}

虽然\co{rcu_assign_pointer()}和\co{rcu_dereference()}
理论上可以用于构建任何可以想象的受RCU保护的数据结构，
但在实践中使用更高层的构造通常更好。
因此，\co{rcu_assign_pointer()}和\co{rcu_dereference()}
原语已被嵌入Linux链表操作API的特殊RCU变体中。
Linux有四种双向链表变体：环形的
\co{struct list_head}和线性的
\co{struct hlist_head}/\co{struct hlist_node}、
\co{struct hlist_nulls_head}/\co{struct hlist_nulls_node}和
\co{struct hlist_bl_head}/\co{struct hlist_bl_node}
对。
前者的布局如
\cref{fig:defer:Linux Circular Linked List (list)}所示，
其中绿色（最左边）方框代表链表头，
蓝色（右边三个）方框代表链表中的元素。
这种表示法很繁琐，因此将如
\cref{fig:defer:Linux Linked List Abbreviated}所示简化，
只显示非头部的（蓝色）元素。

\begin{figure}
\centering
\resizebox{3in}{!}{\includegraphics{defer/Linux_hlist}}
\caption{Linux线性链表（\tco{hlist}）}
\label{fig:defer:Linux Linear Linked List (hlist)}
\end{figure}

Linux的\co{hlist}\footnote{
	``h''代表哈希表，与Linux的双指针环形链表
	相比，它将内存使用减半。}
是线性链表，这意味着头部只需一个指针，
而不是环形链表所需的两个，如
\cref{fig:defer:Linux Linear Linked List (hlist)}所示。
因此，使用\co{hlist}可以将大型哈希表的
哈希桶数组的内存消耗减半。
与前面一样，这种表示法很繁琐，
所以\co{hlist}结构将以与\co{list_head}风格链表
相同的方式简化，如
\cref{fig:defer:Linux Linked List Abbreviated}所示。

Linux的\co{hlist}的一个变体名为\co{hlist_nulls}，
提供多个不同的\co{NULL}指针，
但使用与\cref{fig:defer:Linux Linear Linked List (hlist)}
相同的布局。
在此变体中，低位为零的\co{->next}指针
被认为是一个指针。
然而，如果低位被设为一，则高位标识
\co{NULL}指针的类型。
此类型的链表用于允许无锁（lock-free）读者检测
节点何时从一个链表移动到另一个链表。
例如，哈希表的每个桶可以使用其索引来标记
其\co{NULL}指针。
如果读者遇到与其起始桶索引不匹配的\co{NULL}指针，
该读者就知道它正在遍历的元素在遍历期间
被移动到了其他桶，带着该读者一起。
读者可以使用\apik{is_a_nulls()}函数（当传入
\co{hlist_nulls}的\co{NULL}指针时返回\co{true}）
来确定何时到达链表末尾，
使用\apik{get_nulls_value()}函数
（返回其参数的\co{NULL}指针标识符）
来获取\co{NULL}指针的类型。
当\co{get_nulls_value()}返回意外值时，
读者可以采取纠正措施，例如从头重新开始遍历。

\QuickQuiz{
	但如果一个\co{hlist_nulls}读者被移到某个其他桶
	然后又移回来了呢？
}\QuickQuizAnswer{
	一种处理方式是始终将节点移动到目标桶的开头，
	确保当读者到达具有匹配\co{NULL}指针的链表末尾时，
	它已经搜索了整个链表。

	当然，如果在具有大量元素的哈希表中
	移动操作太多，读者可能永远无法到达链表末尾。
	在常见情况下避免这一问题的一种方法是
	保持哈希表调优良好，从而保持链表较短。
	一种检测问题并处理它的方法是让读者在
	遍历了某个大量节点后终止搜索，
	获取更新侧锁，然后重做搜索，
	但这可能引入死锁。
	另一种完全避免问题的方法是让读者
	在RCU read-side critical section内搜索，
	并在连续更新之间等待RCU grace period。
	一种折中方案是每$N$次更新等待一个RCU grace period，
	对于某个合适的$N$值。
}\QuickQuizEnd

关于\co{hlist_nulls}的更多信息可在Linux内核源代码树中
找到，\path{rculist_nulls.rst}文件
（旧内核中为\path{rculist_nulls.txt}）中提供了
有用的示例代码。

Linux的\co{hlist}的另一个变体结合了位加锁，
名为\co{hlist_bl}。
此变体使用与\cref{fig:defer:Linux Linear Linked List (hlist)}
相同的布局，但保留头指针（图中的``first''）的
低位来锁定链表。
这种方法也减少了内存使用，因为它允许
将本来需要单独的自旋锁存储与指针本身一起。

\begin{sidewaystable*}[htbp]
\rowcolors{1}{}{lightgray}
\renewcommand*{\arraystretch}{1.3}
\centering
\caption{受RCU保护的链表API}
\label{tab:defer:RCU-Protected List APIs}
\footnotesize
\newlength{\cwa}\newlength{\cwb}\newlength{\cwc}\newlength{\cwd}
\IfNimbusAvail{
  \renewcommand{\ttdefault}{NimbusMonoN}
  \setlength{\cwa}{1.9in}\setlength{\cwb}{2.1in}
  \setlength{\cwc}{1.8in}\setlength{\cwd}{1.6in}
}{
  \setlength{\cwa}{1.95in}\setlength{\cwb}{2.15in}
  \setlength{\cwc}{1.9in}\setlength{\cwd}{1.7in}
}
\ebresizewidthsw{
\begin{tabular}{>{\raggedright\arraybackslash}p{\cwa}
    >{\raggedright\arraybackslash}p{\cwb}
    >{\raggedright\arraybackslash}p{\cwc}
    >{\raggedright\arraybackslash}p{\cwd}}
\toprule
\pmb{\tco{list}}：环形双向链表 &
    \pmb{\tco{hlist}}：线性双向链表 &
	\pmb{\tco{hlist_nulls}}：带标记NULL指针的线性双向链表，最多31位标记 &
	    \pmb{\tco{hlist_bl}}：带位加锁的线性双向链表 \\
\midrule
\multicolumn{4}{l}{{\bf 结构}} \\
\tco{struct list_head} &
    \tco{struct}{\tt ~}\tco{hlist_head} ~~~~~~~~~~~~~~
    \tco{struct}{\tt ~}\tco{hlist_node} &
	\tco{struct}{\tt ~}\tco{hlist_nulls_head}
	\tco{struct}{\tt ~}\tco{hlist_nulls_node} &
	    \tco{struct}{\tt ~}\tco{hlist_bl_head}
	    \tco{struct}{\tt ~}\tco{hlist_bl_node} \\
\multicolumn{4}{l}{{\bf 初始化}} \\
\tco{INIT_LIST_HEAD_RCU()} &
    &
	&
	    \\
\multicolumn{4}{l}{{\bf 完整遍历}} \\
\tco{list_for_each_entry_rcu()}
\tco{list_for_each_entry_lockless()}
\tco{list_for_each_entry_srcu()} &
    \tco{hlist_for_each_entry_rcu()}
    \tco{hlist_for_each_entry_rcu_bh()}
    \tco{hlist_for_each_entry_rcu_notrace()}
    \tco{hlist_for_each_entry_srcu()} &
	\tco{hlist_nulls_for_each_entry_rcu()}
	\tco{hlist_nulls_for_each_entry_safe()} &
	    \tco{hlist_bl_for_each_entry_rcu()} \\
\multicolumn{4}{l}{{\bf 恢复遍历}} \\
\tco{list_for_each_entry_continue_rcu()}
\tco{list_for_each_entry_from_rcu()} &
    \tco{hlist_for_each_entry_continue_rcu()}
    \tco{hlist_for_each_entry_continue_rcu_bh()}
    \tco{hlist_for_each_entry_from_rcu()} &
	&
	    \\
\multicolumn{4}{l}{{\bf 逐步遍历}} \\
\tco{list_entry_rcu()}
\tco{list_entry_lockless()}
\tco{list_first_or_null_rcu()}
\tco{list_next_rcu()}
\tco{list_next_or_null_rcu()}
\tco{list_tail_rcu()} &
    \multicolumn{1}{p{1.2in}}{\tco{hlist_first_rcu()}
			      \tco{hlist_next_rcu()}
			      \tco{hlist_pprev_rcu()}} &
	\tco{hlist_nulls_first_rcu()}
	\tco{hlist_nulls_next_rcu()} &
	    \tco{hlist_bl_first_rcu()} \\
\multicolumn{4}{l}{{\bf 添加}} \\
\multicolumn{1}{p{1.2in}}{\tco{list_add_rcu()}
			  \tco{list_add_tail_rcu()}} &
    \tco{hlist_add_before_rcu()}
    \tco{hlist_add_behind_rcu()}
    \tco{hlist_add_head_rcu()}
    \tco{hlist_add_tail_rcu()} &
	\tco{hlist_nulls_add_head_rcu()}
	\tco{hlist_nulls_add_tail_rcu()}
	\tco{hlist_nulls_add_fake()} &
	    \tco{hlist_bl_add_head_rcu()}
	    \tco{hlist_bl_set_first_rcu()} \\
\multicolumn{4}{l}{{\bf 删除}} \\
\tco{list_del_rcu()} &
    \multicolumn{1}{p{1.2in}}{\tco{hlist_del_rcu()}
			      \tco{hlist_del_init_rcu()}} &
	\tco{hlist_nulls_del_rcu()}
	\tco{hlist_nulls_del_init_rcu()} &
	    \tco{hlist_bl_del_rcu()} \\
\multicolumn{4}{l}{{\bf 替换}} \\
\tco{list_replace_rcu()} &
    \tco{hlist_replace_rcu()} &
	&
	    \\
\multicolumn{4}{l}{{\bf 拼接}} \\
\tco{list_splice_init_rcu()}
\tco{list_splice_tail_init_rcu()} &
    \tco{hlists_swap_heads_rcu()} &
	&
	    \\
\bottomrule
\end{tabular}
}
\end{sidewaystable*}

这些链表变体的API成员在
\cref{tab:defer:RCU-Protected List APIs}中汇总。
更多信息可在Linux内核源代码树的\path{Documentation/RCU}
目录和Linux Weekly News~\cite{PaulEMcKenney2024RCUAPI}中找到。

然而，本节的其余部分将展开讨论
\apik{list_replace_rcu()}的使用，因为正是这个API成员
赋予了RCU其名称。
此API成员用于执行更复杂的更新，其中链表中间的
一个具有多个字段的元素被原子地更新，
使得给定的读者看到的要么是旧的值集合，
要么是新的值集合，而不是两个集合的混合。
例如，链表的每个节点可能有整数字段
\co{->a}、\co{->b}和\co{->c}，
可能需要将给定节点的字段从5、6、7分别
更新为5、2、3。

实现此原子更新的代码非常直接：

\begin{fcvlabel}[ln:defer:Canonical RCU Replacement Example (2nd)]
\begin{VerbatimN}[samepage=true,commandchars=\\\[\],firstnumber=15]
q = kmalloc(sizeof(*p), GFP_KERNEL);	\lnlbl[kmalloc]
*q = *p;				\lnlbl[copy]
q->b = 2;				\lnlbl[update1]
q->c = 3;				\lnlbl[update2]
list_replace_rcu(&p->list, &q->list);	\lnlbl[replace]
synchronize_rcu();			\lnlbl[sync_rcu]
kfree(p);				\lnlbl[kfree]
\end{VerbatimN}
\end{fcvlabel}

\begin{figure}
\centering
\IfEbookSize{
\resizebox{2in}{!}{\includegraphics{defer/RCUReplacement}}
}{
\resizebox{2.7in}{!}{\includegraphics{defer/RCUReplacement}}
}
\caption{链表中的RCU替换}
\label{fig:defer:RCU Replacement in Linked List}
\end{figure}

以下讨论逐步分析此代码，使用
\cref{fig:defer:RCU Replacement in Linked List}来说明
状态变化。
每个元素中的三元组分别代表字段\co{->a}、
\co{->b}和\co{->c}的值。
红色底纹的元素可能被读者引用，
由于读者不直接与更新者同步，
读者可能在整个替换过程中并发运行。
请注意，为清晰起见，反向指针和从尾部到头部的
链接被省略了。

链表的初始状态，包括指针\co{p}，
与删除示例相同，如图的第一行所示。

以下文字描述如何以使给定读者
看到这两个值之一的方式，
将\co{5,6,7}元素替换为\co{5,2,3}。

\begin{fcvref}[ln:defer:Canonical RCU Replacement Example (2nd)]
\Clnref{kmalloc}分配一个替换元素，
导致\cref{fig:defer:RCU Replacement in Linked List}
第二行所示的状态。
此时，没有读者可以持有对新分配元素的引用
（如其绿色底纹所示），且该元素未初始化
（如问号所示）。

\Clnref{copy}将旧元素复制到新元素，导致
\cref{fig:defer:RCU Replacement in Linked List}
第三行所示的状态。
新分配的元素仍然不能被读者引用，但现在已初始化。

\Clnref{update1}将\co{q->b}更新为值``2''，
\clnref{update2}将\co{q->c}更新为值``3''，
如\cref{fig:defer:RCU Replacement in Linked List}
第四行所示。
注意新分配的结构仍然对读者不可访问。

现在，\clnref{replace}执行替换，使新元素
终于对读者可见，因此显示为红色底纹，如
\cref{fig:defer:RCU Replacement in Linked List}
第五行所示。
此时，如下所示，我们有两个版本的链表。
已有的读者可能看到\co{5,6,7}元素（因此
现在显示为黄色底纹），但新读者将看到\co{5,2,3}元素。
但任何给定的读者都保证看到一组值或另一组值，
而不是两组的混合。

在\clnref{sync_rcu}上的\co{synchronize_rcu()}返回后，
一个grace period将已过去，因此所有在
\apik{list_replace_rcu()}之前开始的读取都将已完成。
特别是，任何可能持有对\co{5,6,7}元素引用的读者
都保证已退出其RCU read-side critical section，
因此被禁止继续持有引用。
因此，不再有任何读者持有对旧元素的引用，
如\cref{fig:defer:RCU Replacement in Linked List}
第六行中其绿色底纹所示。
就读者而言，我们回到了链表的单一版本，
但新元素替换了旧元素。

在\clnref{kfree}上的\apik{kfree()}完成后，
链表将如\cref{fig:defer:RCU Replacement in Linked List}
最后一行所示。
\end{fcvref}

尽管RCU以替换用例命名，
Linux内核中绝大多数RCU使用依赖于
简单的独立插入和删除，如
\cref{sec:defer:Maintain Multiple Versions of Recently Updated Objects}中
\cref{fig:defer:Multiple RCU Data-Structure Versions}所示。

\QuickQuiz{
	为什么\cref{tab:defer:RCU-Protected List APIs}
	不包含额外的链表初始化API，如\co{INIT_HLIST_HEAD()}？
}\QuickQuizAnswer{
	因为此API和类似的API不是RCU特有的，
	而是也（有时主要是）用于由其他同步机制保护的
	类似链表。
}\QuickQuizEnd

下一节将介绍帮助开发者调试
其使用RCU的代码的API。

\subsubsection{RCU拥有诊断API}
\label{sec:defer:RCU Has Diagnostic APIs}

\begin{table}
\renewcommand*{\arraystretch}{1.15}
\footnotesize
\centering
\begin{tabular}{ll}
\toprule
类别 &
	原语 \\
\midrule
标记RCU指针 &
	\tco{__rcu} \\
\midrule
调试对象支持 &
	\tco{init_rcu_head()} \\
&	\tco{destroy_rcu_head()} \\
&	\tco{init_rcu_head_on_stack()} \\
&	\tco{destroy_rcu_head_on_stack()} \\
\midrule
停滞警告控制 &
	\tco{rcu_cpu_stall_reset()} \\
\midrule
Callback检查 &
	\tco{rcu_head_init()} \\
&	\tco{rcu_head_after_call_rcu()} \\
\midrule
lockdep支持 &
	\tco{rcu_read_lock_held()} \\
&	\tco{rcu_read_lock_bh_held()} \\
&	\tco{rcu_read_lock_sched_held()} \\
&	\tco{srcu_read_lock_held()} \\
&	\tco{rcu_is_watching()} \\
&	\tco{RCU_LOCKDEP_WARN()} \\
&	\tco{rcu_sleep_check()} \\
\bottomrule
\end{tabular}
\caption{RCU诊断API}
\label{tab:defer:RCU Diagnostic APIs}
\end{table}

\Cref{tab:defer:RCU Diagnostic APIs}
展示了RCU的诊断API。

\co{__rcu}标签标记受RCU保护的指针，例如
\qtco{struct foo __rcu *p;}。
可能被传递给\co{rcu_dereference()}的指针可以被标记，
但持有\co{rcu_dereference()}返回值的指针不应被标记。
在变量、结构字段、函数参数和返回值上
提供这些标记，允许Linux内核的\co{sparse}工具
检测受RCU保护的指针被
纯C语言加载和存储错误访问的情况。

调试对象支持对于作为从Linux内核内存分配器获得的
结构一部分的任何\apik{rcu_head}结构是自动的，
但构建自己的专用内存分配器的人
可以分别在分配和释放时使用\apik{init_rcu_head()}和
\apik{destroy_rcu_head()}。
使用在函数调用栈上分配的\co{rcu_head}结构的人
（确实有这种情况！）可以在首次使用前使用
\apik{init_rcu_head_on_stack()}，在最后使用后、
函数返回前使用\apik{destroy_rcu_head_on_stack()}。
调试对象支持允许检测将同一\co{rcu_head}结构
快速连续传递给\co{call_rcu()}及其相关函数的错误，
这是\co{call_rcu()}对应于著名的双重释放类
内存分配错误的等价物。

停滞警告控制由\apik{rcu_cpu_stall_reset()}提供，
它允许调用者在当前grace period的剩余时间内
抑制RCU CPU停滞警告。
RCU CPU停滞警告有助于定位
RCU read-side critical section运行时间过长的情况，
内核调试器等可以抑制它们是有用的，
例如在遇到断点时。

Callback检查由\apik{rcu_head_init()}和
\apik{rcu_head_after_call_rcu()}提供。
前者在将\co{rcu_head}结构传递给\co{call_rcu()}之前
对其调用，然后\co{rcu_head_after_call_rcu()}将
检查callback是否已用指定函数调用。

对lockdep~\cite{JonathanCorbet2006lockdep}的支持包括
\apik{rcu_read_lock_held()}、
\apik{rcu_read_lock_bh_held()}、
\apik{rcu_read_lock_sched_held()}和
\apik{srcu_read_lock_held()}，
每个在相应类型的RCU read-side critical section内
调用时返回\co{true}。

\QuickQuiz{
	为什么没有用于Tasks RCU的\co{rcu_read_lock_tasks_held()}？
}\QuickQuizAnswer{
	因为Tasks RCU没有读侧标记。
	相反，Tasks RCU的read-side critical section
	由自愿上下文切换界定。
}\QuickQuizEnd

因为\co{rcu_read_lock()}不能在空闲循环中使用，
且能效考虑已导致空闲循环变得相当复杂，
\apik{rcu_is_watching()}在可以合法使用\co{rcu_read_lock()}
的上下文中返回\co{true}。
再次注意\co{srcu_read_lock()}可以在空闲甚至
离线的CPU上使用，这意味着\co{rcu_is_watching()}
不适用于SRCU。

\apik{RCU_LOCKDEP_WARN()}在lockdep启用且
其参数求值为\co{true}时发出警告。
例如，\co{RCU_LOCKDEP_WARN(!rcu_read_lock_held())}
在RCU read-side critical section之外调用时会发出警告。

\apik{RCU_NONIDLE()}可用于强制RCU在执行
作为唯一参数传入的语句时保持监视。
例如，\co{RCU_NONIDLE(WARN_ON(!rcu_is_watching()))}
永远不会发出警告。
然而，2020--2021年期间的变更将RCU的覆盖范围
扩展到空闲循环的更深处，这应该大大减少甚至
消除对\co{RCU_NONIDLE()}的需要。

最后，\apik{rcu_sleep_check()}在RCU、RCU-bh或
RCU-sched read-side critical section内调用时发出警告。

\subsubsection{RCU的API可以在哪里使用？}
\label{sec:defer:Where Can RCU's APIs Be Used?}

\begin{figure}
\centering
\resizebox{3in}{!}{\includegraphics{defer/RCUenvAPI}}
\caption{RCU API使用约束}
\label{fig:defer:RCU API Usage Constraints}
\end{figure}

\Cref{fig:defer:RCU API Usage Constraints}
展示了哪些API可以在哪些内核内环境中使用。
RCU读侧原语可以在任何环境中使用，包括NMI；
RCU变更和异步grace period原语可以在
NMI之外的任何环境中使用；
最后，RCU同步grace period原语只能在进程上下文中使用。
RCU链表遍历原语包括\apik{list_for_each_entry_rcu()}、
\apik{hlist_for_each_entry_rcu()}等。
类似地，RCU链表变更原语包括
\apik{list_add_rcu()}、\apik{hlist_del_rcu()}等。

请注意，作为一般规则，其他RCU系列的原语可以替代使用，
不过\co{srcu_read_lock_nmisafe()}而非
\co{srcu_read_lock()}可以在NMI上下文中使用。

\subsubsection{那么，RCU\emph{到底}是什么？}
\label{sec:defer:So; What is RCU Really?}

在其核心，RCU不多不少就是一个支持
插入的发布和订阅、等待所有RCU读者完成、
以及维护多个版本的API。
话虽如此，可以在RCU之上构建更高层的构造，
包括\cref{sec:defer:RCU Usage}中列出的
读写锁、引用计数和存在性保证构造。
此外，我毫不怀疑Linux社区将继续
发现RCU的有趣新用途，
正如他们对内核中许多同步原语所做的那样。

当然，对RCU更完整的理解还应包括
这些API可以做的所有事情。

然而，对于许多人来说，对RCU的完整理解
必须包括RCU实现的示例。
\Cref{chp:app:``Toy'' RCU Implementations}因此展示了一系列
复杂度和能力递增的"玩具"RCU实现，
不过有些人可能更喜欢经典的
``User-Level Implementations of Read-Copy
Update''~\cite{MathieuDesnoyers2012URCU}。
对于其他所有人，下一节给出了一些RCU用例的概述。
